\section{Conclusions}
\label{sec:conclusions}

This review analysed 31 studies, published from 2021 onwards, that prioritise industrial maintenance interventions and evaluate their approach in a real setting. The building blocks are mature: criteria centred on resources, failure likelihood and production impact, weighting and ranking methods such as AHP, BWM and TOPSIS, optimisation models for crews, budgets and time windows, and increasingly reliable failure predictions. What remains open is how these blocks are put together. Safety is seldom treated as a separate condition, ranking and scheduling are usually solved apart, priorities are rarely updated, weights have a weak basis, recommendations are only partly explained, and few studies show real improvements in maintenance results.

The proposed DSS addresses part of these gaps by combining failure prediction with AHP and TOPSIS, placing safety outside the trade-off between criteria and explaining every position in the ranking. The next steps are to implement the prototype on the Microsoft Azure dataset characterised in Section~\ref{sec:dataset} and Appendix~\ref{app:eda}, elicit the weights with maintenance, production and safety staff, and evaluate it in scenarios, comparing its ranking with that of an experienced planner, testing its sensitivity to the weights and measuring whether users accept the explanations. Integration with scheduling and validation with real plant data are left for later stages.

In our view, as digital twins and condition monitoring make more data available, the value of DSS in industrial maintenance will lie less in replacing the planner than in making decisions explicit: showing why an intervention comes first, helping departments agree on that order and keeping a record from which the plant can learn.